\chapter{知识的演进与迁移概述}
\label{chap:ket-overview}

一个花卉识别系统已经能够根据照片判断物种。把它交给另一地区的使用者，原有判断还能否成立？要求它识别叶片病害时，过去学到的形状与纹理是否能减少标注需求？如果每个月都加入新物种，怎样避免新版本只认识最近见过的花卉？这些问题都以已经学得的经验为起点，但要求保留的能力、需要改变的对象和判断成功的标准并不相同。

按照全书的内容安排，答案、关系、约束与目标的获取属于第~\ref{part:efficient-knowledge-use}部分，交互决策与系统实现分别属于第~\ref{part:reinforcement-learning}、\ref{part:systems}部分。这三部分在当前版本中尚未展开，本部分不以其正文作为前置阅读。这里从已经形成的表示、参数、预测行为和记忆出发，考察它们怎样继续支持学习与使用。理解这一问题不需要预先掌握某种迁移算法，却需要把四件事分开：实际需求发生了什么变化，已有信息保存在何处，当前可以访问和改变什么，以及改变后的收益应由什么证据确认。本章据此建立共同概念，再说明后续研究问题及其历史联系。

本部分采用教程式综述而非穷尽性系统综述，文献覆盖截至2026年10月。选取的工作须至少承担一种明确职责：奠定概念或问题边界，代表一条仍有解释力的机制路线，提供独立比较或负面结果，或反映2022年以后改变方法地图的进展。正文优先展开可复算机制、适用条件和失败边界，不按论文数量推断路线重要性，也不把提出论文中的结果视为跨任务稳定结论。数据获取、强化学习决策、具体系统工程以及可信性保证只在与知识演进接口直接相关时出现，其完整方法留给相应部分；近期预印本只用于标示证据方向，不单独支撑确定性结论。

\section{知识的演进与迁移的场景}
\label{sec:ket-overview-scenarios}

设想花卉系统最初使用某地春季拍摄的照片训练。新地区仍要识别相同物种，但背景、光照、常见品种的比例和拍摄距离都可能不同。原模型辨认花瓣边缘的经验也许仍然有用，依赖某种背景颜色的判断却可能失效。进入另一季节时，同一物种还可能呈现不同生长状态。这里不能仅凭“识别对象没有变”就断定模型无需调整，也不能仅凭照片外观改变就认定所有经验都已过时。

需求也可能改变在同一照片上要回答的问题。物种识别问的是“这是什么花”，病害判断问的是“叶片是否出现某类异常”，花朵定位还要求指出对象所在的位置。过去形成的局部纹理与轮廓表示可以成为候选起点，但物种名称不能充当病害标签，整张图的类别分数也不能直接充当位置坐标。若只得到少量新物种照片，则要检查旧表示是否保留了区分这些物种所需的细节。新地区同时出现新物种时，环境变化与目标类别扩展会一起发生。

当系统需要同时输出物种、病害和位置时，让多个任务共同学习可能减少重复计算，并让一种标注补充另一种标注没有强调的信息。不过，共同使用一套表示也会引入竞争：物种分类希望忽略的斑点，恰好可能是病害判断所依赖的证据。若这些任务不是同时开放训练，而是按月份逐个加入，系统还必须在吸收新经验时保持已承诺的旧能力。只检查最新任务，无法发现原来会识别的物种是否已经被新更新破坏。

更换传感器揭示了另一种差别。把普通彩色相机换成多光谱传感器，仍可要求输出同样的物种名称，但输入通道的物理含义已经不同。即使把两类输入整理成相同长度的向量，也不能认为对应坐标包含相同信息。系统需要建立新观测与旧表示之间有根据的联系，而不是仅仅延长训练时间。另一方面，若两个团队已经分别训练了擅长不同地区或不同病害的模型，直接需求是让现成能力共同发挥作用；把原始数据集中起来从头训练只是一个可能的对照，并不等于已经整合了这两个模型。

变化还可以来自知识本身。假设系统附带植物研究机构的问答服务，机构负责人在明确日期由甲变为乙，当前问答应随之修改，历史日期的答案却应保留甲。触发更新的是一条经过核对的事实请求，不必等到照片分布发生变化。从头训练可能重复已有工作，完全不改又会继续给出过时回答。我们需要确定修改作用于外部记录、模型参数还是本次输入，并检查它影响的范围。

图~\ref{fig:ket-overview-scenarios}把这些请求放在同一个已训练系统周围。各分支描述主要变化对象，不规定执行顺序；一次上线可以同时遇到其中几项。决定后续做法之前，应能指出哪些经验仍有使用依据，以及哪些预测要求、观测关系或事实有效范围必须重新确定。

\begin{figure*}[htbp]
  \centering
  % Conceptual relationships, not a temporal pipeline or measured transfer gains.
% Final width is below 169 mm. All outgoing requests share the east-side port.
\begin{tikzpicture}[
  x=1mm,y=1mm,
  font=\figurefont\small,
  every node/.style={text=BookInk},
  source/.style={
    rectangle,draw=BookTeal,fill=BookTeal!8,
    line width=1pt,text width=39mm,minimum height=22mm,
    inner sep=2mm,align=center
  },
  request/.style={
    rectangle,draw=BookBlue,fill=BookPaper,
    line width=0.8pt,text width=51mm,minimum height=12mm,
    inner sep=1.5mm,align=center
  },
  link/.style={draw=BookTeal,line width=0.9pt},
  arrow/.style={link,-{Latex[length=2mm,width=1.4mm]}},
  change/.style={font=\figurefont\footnotesize,anchor=south,text=BookMuted}
]
  \node[source] (system) at (23,0) {已训练的\\花卉识别系统};
  \node[font=\figurefont\footnotesize,align=center,text width=42mm]
    at (23,-18) {已有表示、模型与记忆\\附带机构问答服务};

  \node[request] (environment) at (122,36) {新环境适应\\地区、季节与传感器变化};
  \node[request] (target) at (122,18) {新任务学习\\病害、定位与新物种};
  \node[request] (sharing) at (122,0) {多任务共享\\多个目标共同学习};
  \node[request] (editing) at (122,-18) {事实修正\\负责人交接与有效日期};
  \node[request] (integration) at (122,-36) {多来源整合\\接入分别训练的其他模型};

  \coordinate (fork) at (61,0);
  \draw[link] (system.east) -- (fork);
  \draw[link] (61,-36) -- (61,36);
  \foreach \target/\height/\object in {
    environment/36/数据与观测,
    target/18/目标语义,
    sharing/0/任务关系,
    editing/-18/事实有效范围,
    integration/-36/知识来源
  } {
    \draw[arrow] (61,\height) -- (\target.west);
    \node[change] at (78,\height+1) {\object};
  }
  \node[font=\figurefont\footnotesize,text=BookMuted,anchor=north]
    at (93,-46) {各类变化可以同时发生};
\end{tikzpicture}

  \caption{花卉识别系统中的知识使用与变化。箭头从已有系统指向不同使用请求，线旁标出主要变化对象；它们不是必须依次完成的阶段。新地区、新物种与传感器更换可以共同出现，事实修正则使用系统附带的机构问答服务作为示例。多来源整合还需要接入分别训练的其他模型。}
  \label{fig:ket-overview-scenarios}
\end{figure*}

\section{知识的演进和迁移的基本概念}
\label{sec:ket-overview-concepts}

同一项请求可以同时被描述为“换了环境”“需要少量标签”或“按时间持续更新”。这些描述分别涉及变化对象、可用信息和经验顺序，不能互相替代。为了比较方法，需要先确定被使用的信息及其载体，再明确来源与目标，随后记录学习关系和访问条件。

\subsection{知识与知识载体}
\label{sec:ket-overview-carriers}

一个识别器训练结束后，能够留下的不只有一组权重。照片可以继续作为参照，特征提取器可以提供表示，预测接口可以告诉另一个模型它怎样判断。本部分把能够支持后续预测与学习的信息统称为\term{知识}（knowledge），把保存或提供这些信息的具体对象称为\term{知识载体}（knowledge carrier）。这里的“知识”是方法讨论中的操作性概念，不意味着模型内部一定存有能逐条读出的语义事实。

样本与\term{外部记忆}（external memory）保留较直接的经验：系统可以重新读取历史照片、查询某个物种的描述，或找出过去任务中的代表案例。样本通常保留输入和标签之间的联系，外部记忆还可保存出处、时间与适用条件。存储完整照片与只存一个类别名称不同，后者不能重新提供照片中的局部纹理；保存一条事实与保证它被正确检索、正确使用也不同。载体中有什么信息，以及调用过程是否能取到它，需要分别检查。

特征表示把经验转化为供后续计算使用的坐标。记输入为$\bm{x}$，参数为$\bm{\theta}$的特征提取器为$g_{\bm{\theta}}$，输出$\bm{z}=g_{\bm{\theta}}(\bm{x})\in\mathbb R^d$，其中$d$是特征维数。冻结$g_{\bm{\theta}}$并训练新的病害预测头，复用的是已有映射所保留的区分信息；目标训练不改变该映射。关系表示则可以保存对象之间的联系，例如某个病害与哪些可观察症状有关。无论采用向量还是显式关系，表示都可能丢失后续任务需要的信息，不能仅以它曾经支持过一个任务就断言它适合所有目标。

模型参数与模块提供另一种复用方式。把物种模型的参数作为病害模型的初始化，允许后续训练调整原表示，使用的是一个学习起点，而非永久固定的特征坐标。复用局部模块时，还要保留它与其他模块相接的输入输出含义。相比之下，只能查询的教师模型通过\term{预测行为}（predictive behavior）提供信息，也就是它对给定输入实际给出的类别、分数或概率。另一个模型可以学习这些反应，却未必能够取得教师的中间表示或重建全部训练经验。

经验也可以改变以后怎样学习。历史任务可以帮助选择\term{参数先验}（parameter prior），即在当前样本不足时偏好的参数范围及相对可能性；也可以帮助形成\term{学习规则}（learning rule），即根据新证据决定下一步如何修改模型的规则。前者约束倾向采用什么解，后者约束怎样到达解。按任务保存的少量案例还可以构成任务记忆，在遇到相似目标时提供参照或选择依据。冻结表示、参数初始化、教师输出与任务记忆因而是同一段经验发挥作用的不同途径，访问要求和失效方式并不相同。

\begin{example}[带生效日期的事实修正]
\label{ex:ket-overview-dated-fact}
以下机构、人物和日期均为教学构造。假设某植物研究机构在2026年7月1日由负责人甲交接给乙，此前由甲任职。系统在8月被问“该机构现在的负责人是谁”时应回答乙，被问“2026年6月的负责人是谁”时仍应回答甲。

修改外部记录，可以把甲、乙分别与各自有效日期关联；这改变了可检索证据，尚需检查问答过程是否读取了相应版本。修改模型参数，改变的是模型从问题到答案的计算，需检查不同问法和日期条件下的输出，而不能只确认参数文件已保存。仅在当前问题前补充交接说明，改变的是这次预测的条件；若不另行保存或训练，不能据此声称下次独立查询也会使用新事实。
\end{example}

示例~\ref{ex:ket-overview-dated-fact}说明，可观察的\term{任务能力}（task capability）与内部存储是不同层次：前者是系统在规定输入、工具和评价条件下能完成什么，后者是模型或记忆保存了什么。参数发生变化不能单独证明某条事实已经写入，检索返回了正确文本也不能单独证明最终回答使用了它。反过来，一次回答正确也不足以定位这项能力究竟来自参数、检索还是当前上下文。更新对象及其验证标准将在第~\ref{chap:ket-update-maintenance}章进一步对应起来。

\subsection{来源与目标}
\label{sec:ket-overview-source-target}

判断一段经验有没有用，需要说明它将被用于什么。\term{来源}（source）提供已有经验、模型或学习规则，\term{目标}（target）规定当前希望完成的预测与学习要求。两者是相对于当前请求的角色：今天形成的目标模型，明天可以成为另一个请求的来源。来源也不必是一份仍然可下载的数据集，一个只能调用的模型接口同样能提供来源信息。

\term{领域}（domain）主要描述数据的观测空间及其分布，也就是会看到什么样的输入、这些输入怎样出现；\term{任务}（task）还要说明输出表示什么，以及怎样判断回答是否成功。对于分类任务，只写“输出一个整数”不够，必须知道整数对应哪个物种或病害；只写“准确率较高”也不够，还要确定在哪些对象和哪种错误计分规则下评价。领域与任务在文献中有不同形式化约定，本书始终分别记录预测要求、数据分布与经验到达方式，以避免名称掩盖实质条件。

以可用监督目标表示的预测问题为例，记来源和目标的输入空间为$X_S,X_T$，输出空间为$Y_S,Y_T$，数据集为$S_S,S_T$。联合分布$\mathcal P_S,\mathcal P_T$分别描述相应输入与标签共同出现的规律。设$f_D:X_D\to Y_D$为一个预测器，$\ell_D(\hat y,y)$为任务损失，其中$D\in\{S,T\}$只用来选择来源或目标。风险用来表示这个预测器在相应分布下的平均损失：
\begin{equation}
  R_D(f_D)=
  \mathbb E_{(\bm{x},y)\sim\mathcal P_D}
  \bigl[\ell_D(f_D(\bm{x}),y)\bigr],
  \qquad D\in\{S,T\}.
  \label{eq:ket-overview-risk}
\end{equation}
若两侧输入输出接口与损失相同，可以把同一个$f$放入$R_S(f)$和$R_T(f)$，比较它在两域的表现。若任务从物种分类变为定位，输出结构与损失含义不同，两个风险的数值就不能直接作为同一指标相减。目标风险也不是总能由现有数据算出：没有目标标签时，观察到目标输入并不等于观察到预测错误。

对于含$n_T$个带标签样本的目标数据集$S_T=\{(\bm{x}_i,y_i)\}_{i=1}^{n_T}$，相应经验风险为
\begin{equation}
  \hat R_T(f_T)=\frac{1}{n_T}
  \sum_{i=1}^{n_T}\ell_T(f_T(\bm{x}_i),y_i).
  \label{eq:ket-overview-empirical-risk}
\end{equation}
下标$i$标记样本，帽号区分有限样本上的平均与分布上的期望。算法可以用经验风险构造训练目标，但当前要验证的仍是未来目标表现。若还有旧能力保留或存储预算要求，它们需要进入约束或附加目标，不能指望最小化当前任务损失自动满足。

同类花卉换拍摄地区时，通常保持$X_S=X_T=X$和$Y_S=Y_T=Y$，也保持物种标签含义与损失，变化的是$\mathcal P_S$与$\mathcal P_T$。输入边缘分布$\mathcal P_D^X$描述各种照片出现的比例，条件关系$\mathcal P_D(y\mid\bm{x})$描述观察到一张照片后标签的分布。地区改变可能影响前者；即使标签语义固定，照片未能记录的生长条件或被遮挡的辨识特征也可能使后者改变。不能从场景名称自动推出条件关系不变。只有额外规定同一输入的标签规律保持时，才进入理论卷所述的协变量偏移设定。若改变的是物种划分或标注标准，则应重新核对任务定义，不能把它混作固定任务下的环境变化。

同一照片增加病害标签时，输入可以完全不变，输出含义却由物种变为病害。即使两种标签碰巧都编码为$1,\ldots,c$，整数相同也不能建立语义对应。仍然识别同类花卉但更换传感器时，则可能保留$Y_S=Y_T$与损失，却有$X_S\neq X_T$；此时应问新观测能否提供旧预测所需的证据，而不是直接比较两种输入坐标。联合分布、条件关系与任务损失由此分别描述数据、预测规律和成功标准。

经验怎样到达是另一项记录。一次性收到全部目标照片，与按季节收到$S^{(1)},S^{(2)},\ldots$，即使最终包含相同样本，也可能允许不同的训练和保存操作；带括号的上标表示到达或学习阶段。理论卷第~\ref{chap:transfer-learning-theory}章给出分布联系、风险界与可识别性的完整讨论，第~\ref{sec:distribution-shift}节区分具体偏移。这里特别保留定理~\ref{thm:transfer-nonidentifiability}的提醒：缺少必要联系或覆盖时，无标签目标输入无法补出任意未知的标签规律。

\subsection{演进与迁移的学习设定}
\label{sec:ket-overview-learning-settings}

明确来源和目标后，还要区分在考察哪一种关系。\term{知识复用}（knowledge reuse）表示已有信息参与当前预测或学习，例如把保存的照片再次作为参照；它不要求任务发生变化。\term{知识迁移}（knowledge transfer）强调来源经验对另一任务或环境的作用，例如物种表示是否帮助病害识别。\term{知识演进}（knowledge evolution）强调信息随经验积累、重组、修正或遗忘的过程，例如系统逐渐增加物种，同时撤下失效事实。演进不等于训练步数增多：对同一数据反复优化可以只是在求解一个固定问题，尚未回答新知识如何进入、旧知识为何保留。

复用也不总是持久写入参数。把少量新物种示例放进当前输入，模型可能据此完成一次条件预测；撤去这些示例后，行为可能恢复原状。这属于利用\term{上下文}（context），即本次预测可读取的附加任务证据。若把示例保存进任务记忆，未来仍可重新读取；若再用它们训练参数，还可能形成不依赖本次输入的变化。因此，临时条件化、持久外部记忆与持久参数修改能够结合，却必须分别记录，不能由一次条件预测宣称完成了长期学习。

本书沿用\term{迁移学习}（transfer learning）描述利用来源知识帮助目标问题的学习设定。“帮助”是研究目标，并非使用了来源就已经实现的结果：目标效果可能变差，或者效果相近但学习代价降低。\term{多任务学习}（multi-task learning, MTL）考察多个目标共同学习时怎样共享信息，例如用物种和病害的训练信号共同塑造表示。其关键是目标之间通过学习过程发生联系，不要求每一个计算步同时包含所有任务，也不要求所有任务使用同一组照片。

当希望系统对尚未到来的任务学得更好时，训练对象还可以包括学习起点或学习规则。\term{元学习}（meta-learning）利用多个任务的经验改善后续学习，也就是让准备阶段为“之后怎样学”服务，而不只为当前任务求一个预测器。用$\mathcal T_k$表示第$k$个任务，其中规定输入、输出语义、损失及任务内分布；用$\Pi$表示这些任务的抽样分布。写$\mathcal T_k\sim\Pi$，指抽取一个学习问题，而不是从某个固定问题中抽一张照片。能在训练任务上拟合良好，与能利用少量经验适应从$\Pi$抽出的新任务，属于不同要求；超出这个任务分布的目标还需要额外检查。

若经验受到时间顺序限制，系统还要在适应新情况时保留约定能力，这构成\term{持续学习}（continual learning）。它关心多次学习之间的关系，而非只关心某次适应是否完成。连续经验可以来自环境变化，可以不断增加类别，也可以引入新的预测任务。旧照片未必能够无限保存，任务边界也未必可见；“把所有历史数据重新混在一起训练”因而不总是允许的解决方案。

考虑一个完整过程：先在大量花卉照片上预训练表示，再用物种与病害标签共同微调，之后依次接收新的定位任务和新增物种。来源表示参与目标训练体现迁移，多个目标共同影响表示体现多任务共享，新任务逐次进入并要求旧能力不退化体现持续学习。如果准备阶段还明确优化利用少量任务经验后的适应效果，就进一步具有元学习目标；这种适应可以依赖参数更新，也可以依赖前向计算对任务上下文的使用。仅仅先训练再微调，并不能证明进行了元学习。这些设定刻画不同关系，不是按训练先后、任务数量或标签多少划成的一组互斥类别。

\subsection{知识复用的信息条件}
\label{sec:ket-overview-information}

两个方法即使都被称为迁移学习，也可能没有使用同样的信息。源侧能够访问完整数据时，可以重新取样、检查覆盖或回放历史经验；能够访问模型参数时，可以复制、冻结或调整指定部分。读取中间表示、取得某个可执行模块和只查询预测接口又是不同权限：黑盒接口返回的类别概率不会直接提供中间特征，也不能直接提供对源参数的梯度。相应地，源数据不可访问只排除了直接使用原始样本，并不意味着已经训练好的模型不再携带来源信息。

目标侧应分别记录标签、无标签输入、类别或任务描述，以及反馈何时可用。只有输入时，系统可以估计某些输入统计，却不能直接计算式~\eqref{eq:ket-overview-empirical-risk}中的监督损失。任务描述规定输出含义，有时还提供可复用的语义联系，但不能一律替代真实标签。预测之后才到达的病害鉴定结果，可以用于以后更新；若评价要求先预测后反馈，就不能把该结果用于已经发生的那次预测。即使标签总数相同，反馈时机不同也会改变问题难度。

表~\ref{tab:ket-overview-information}采用同一个教学目标：在新地区识别原来的十种花卉，同时检查旧地区能力。协议A允许离线访问源数据和源模型，并提供20张带标签目标照片与200张无标签目标照片；协议B只允许查询冻结源模型的十类概率，200张无标签目标照片按拍摄时间逐张到达，不提供目标标签或地区身份。A可修改本地模型，B仅可更新协议允许的外部状态或另建承接模型，不得修改源服务；两者的目标类别含义相同。表中数量仅为说明信息条件的构造，不是实验结果。

\begin{table*}[htbp]
  \centering
  \small
  \begin{tabularx}{\linewidth}{@{}p{17mm}XXXX@{}}
    \toprule
    信息对象 & 可用时机 & 访问方式 & 允许用途 & 受限时影响的操作 \\
    \midrule
    源数据 & A：离线准备和更新；B：不可用 & A：本地带标签照片；B：不开放 & A：检查覆盖、回放及监督训练 & B不能直接重取源样本或重新标注 \\
    源参数 & A：全程；B：不可用 & A：读写权重及中间表示；B：不开放 & A：冻结特征、初始化或局部更新 & B不能微调源模型或按源参数融合 \\
    源预测 & A：本地随时；B：当前输入到达后 & A：本地前向；B：查询十类概率 & 均可在允许的输入上取得教师行为 & B不提供中间量，且受接口次数限制 \\
    目标标签 & A：适应前20张；B：适应中不可用 & A：读取标签；B：仅独立评测方持有 & A：监督适应与选择，分别记账 & B不能用真实标签选模型或算监督损失 \\
    目标无标签数据 & A：适应前200张；B：逐张到达 & A：可重复遍历；B：只读当前及已存历史 & A：整批统计；B：因果使用已见输入 & B不能提前读取未来输入统计 \\
    任务身份 & A、B：固定物种任务均已知；地区身份仅A可见 & 均知任务含义；A另读地区元数据 & A可据地区标记选择模块 & B不能借地区身份路由；任务与地区须分开记录 \\
    时间顺序 & A：离线完整批次；B：严格在线顺序 & A：可重排；B：按到达次序 & A可批量适应；B只用当时已有信息 & 未来样本或延迟反馈不得提前使用 \\
    历史状态 & A：可保存样本和检查点；B：仅有界缓存 & A：本地读写；B：按预定容量保存 & A可回放和回退；B可复用保留状态 & B须计入淘汰损失，不能假设全史可得 \\
    \bottomrule
  \end{tabularx}
  \caption{同一花卉目标任务下的两套信息条件。A与B的具体含义见正文，表中的不可用均指学习或选择阶段的权限，不排除独立评测方持有测试标签。两套协议都可能被宽泛地称为跨域迁移，但监督、源接口、时间顺序与历史访问不同，效果差异不能全部归于算法。实际实验还须给定查询次数、缓存容量、计算预算和旧能力保留标准。}
  \label{tab:ket-overview-information}
\end{table*}

在这些条件上，常用设定才有可检查的含义。\term{领域自适应}（也称领域适应，domain adaptation）允许学习过程使用目标域样本，借来源信息改善该目标域上的预测。是否有目标标签、是否能访问源数据，需要另行限定。\term{领域泛化}（domain generalization）按本书采用的标准协议，在训练与模型选择阶段都不使用目标域数据，要求对未见领域作出预测。预测时读取当前输入不等于得到针对目标域调参的权限，这一区别与第~\ref{sec:domain-generalization}节一致。

\term{测试时适应}（test-time adaptation, TTA）允许系统在使用阶段利用规定的信息改变预测所依赖的状态，例如利用当前无标签批次调整部分参数或统计量。协议必须进一步规定能否看完整批次、能否跨样本保留状态、何时重置，以及是否允许反馈。若部署时又利用目标样本调整模型，就不能仍把全过程报告为不使用目标数据的标准领域泛化。表中的B保留了严格顺序约束，但“在线”本身并不规定具体适应算法或保证适应有效。

持续学习也需要按预测时的信息区分。标准\term{任务增量学习}（task-incremental learning）在测试时提供任务身份，例如先告知要在某组物种中分类，系统可据此选择相应输出头；这项身份不同于待预测的类别。\term{类别增量学习}（class-incremental learning）逐步引入新类别，测试时须在截至当前的累积类别中判断，不能靠一个给定标识把旧类排除。\term{领域增量学习}（domain-incremental learning）保持输出类别与任务含义，经历不同输入领域，并在标准比较中不依赖测试时给出的领域身份。文献协议可能采用不同约定，因此任务身份在训练、选择和预测时是否可见，都应显式说明。

任务身份、历史访问与允许更新对象之外，还必须写出资源预算和保留要求。固定缓存容量与每个任务固定容量不是同一长期条件；必须继续服务旧地区，与可以只优化新地区也不是同一目标。授权使用、隐私泄漏与恶意来源的风险分别需要可信性部分的规范和证据，相关讨论见第~\ref{chap:trust-privacy}、\ref{chap:trust-security}章。只开放预测接口、隐藏源数据或把数据留在原地，都不自动构成隐私或安全保证。

\section{知识的演进与迁移的研究内容}
\label{sec:ket-overview-research}

面对一个新请求，直接选取某个算法往往太早。花卉系统进入新地区之前，需要判断来源经验相对于无来源起点是否有收益，以及收益体现为预测更准确、所需标签更少还是总成本降低。即使决定复用，也要选择具体样本、网络层、模块或来源模型；整个模型在旧地区准确，不等于每个部分都适合当前目标。这些判断组成第~\ref{chap:ket-reuse-selection}章的三个相连问题，评价条件采用上一节建立的信息记录。

作出选择后，研究目标决定接下来要解释什么。如果仍服务于已有系统，但环境、观测或知识内容发生变化，中心问题是更新与维护：吸收新证据，保持仍然有效的能力，修正错误事实，并在必要时消除指定训练记录的影响。保留当前成绩还不够，长期系统也要维持\term{可塑性}（plasticity），即继续从以后经验中有效学习的能力。事实撤回与普通识别错误不同，不能把系统不再输出某个答案直接当成训练影响已消除；删除保证及其验证要求见第~\ref{sec:privacy-unlearning}节。

如果要形成病害识别、定位或新物种辨认能力，重点转为来源经验怎样支持新的输出语义，以及达到目标水平花费了多少学习代价。这是迁移与适应的问题。新任务连续进入时，它又与维护要求相接：学会本月新物种是一项成功，仍能在旧新物种的共同集合中可靠分类是另一项成功。前者由第~\ref{chap:ket-transfer-adaptation}章说明目标能力的形成，后者由第~\ref{chap:ket-update-maintenance}章说明累积能力的保留；环境或传感器差异所需的基础校正机制也在更新章集中建立。

多个任务共同参与训练时，需要决定共享哪些计算、怎样协调目标以及如何分配样本和更新机会；这是第~\ref{chap:ket-sharing}章的起点。多个模型已经分别训练完成时，则需要决定保留它们共同计算，组合其参数，还是训练另一个模型承接其行为；这是第~\ref{chap:ket-integration}章的起点。整合中仍然可能训练组合器或承接模型，因此不能用“有没有训练”区分整合与共享，区别在于当前要利用的是现成来源能力，还是多个目标共同形成知识的过程。

表~\ref{tab:ket-overview-chapters}汇总这些研究问题及其组织方式。它是一张问题地图，而不是唯一的处理流水线；评价也不是等所有更新结束后才补做的一步。必须在选择和修改之前明确成功标准，在过程中保留必要记录，才能由第~\ref{chap:ket-evaluation}章的评价框架判断实际收益。

\begin{table*}[htbp]
  \centering
  \small
  \begin{tabularx}{\linewidth}{@{}p{47mm}XX@{}}
    \toprule
    章节 & 核心研究问题 & 主体节的组织依据 \\
    \midrule
    第~\ref{chap:ket-overview}章：知识的演进与迁移概述 & 为什么需要演进与迁移，研究什么，怎样发展至今 & 场景、基本概念、研究内容与研究历史 \\
    第~\ref{chap:ket-reuse-selection}章：知识复用的判断与选择 & 复用是否有收益，哪些知识可复用，应当选择哪些 & 收益判断、可复用知识判断与知识选择 \\
    第~\ref{chap:ket-update-maintenance}章：知识的更新和维护 & 怎样改变已有知识并维持系统能力 & 确定目标、准备依据、更新、验证与持续维护 \\
    第~\ref{chap:ket-integration}章：知识的整合 & 怎样整合分别学得的知识，形成所需组合能力 & 模型组合、参数融合与知识蒸馏 \\
    第~\ref{chap:ket-transfer-adaptation}章：知识的迁移与适应 & 怎样利用来源经验形成新任务能力 & 明确目标、准备可迁移知识与适应目标 \\
    第~\ref{chap:ket-sharing}章：知识的共享 & 多个任务共同学习时怎样协调共享 & 任务组织、共享设计、联合目标与训练资源 \\
    第~\ref{chap:ket-evaluation}章：知识演进与迁移的评价 & 怎样证明变化产生了预期收益 & 目标指标、数据基准、实验协议与证据分析 \\
    \bottomrule
  \end{tabularx}
  \caption{本部分的研究问题与章节关系。章节按直接研究目标组织，各章内部可以采用过程、机制或设计决策的划分方式。同一请求可以调用多章，但某项基础机制只在主要归属处完整展开。}
  \label{tab:ket-overview-chapters}
\end{table*}

这张章节地图与后面的历史四线使用同一组坐标：共享规律与快速适应主要解释来源经验为何可能支持复用、迁移和共享；顺序学习解释更新与维护为何必须同时检查保留和可塑性；多来源整合解释现成模型如何通过计算、参数或监督进入结果。第~\ref{chap:ket-evaluation}章横切所有线索，统一核对目标、信息权限、时间协议和成本。历史线索回答问题怎样形成，章节职责回答当前决策在哪里展开，两者不是两套并列分类。

回到联合学习多个任务、适应新任务、持续接收新经验的花卉系统。物种与病害任务共同训练后，表示可能更适合新的定位任务，也可能因为丢失空间细节而不适合；有限目标试训可以反过来修正对来源价值的判断。定位适应结束后，若更新过共同表示，还要检查原物种和病害任务。若为保护旧任务而把所有表示完全固定，又可能限制新任务所需的信息。这些反馈说明，来源选择、快速适应、共享和保留需要相互协调，不能只优化其中一个环节的局部指标。

同一种操作也可能服务不同目标。\term{知识蒸馏}（knowledge distillation）让已有模型的预测或中间表示指导另一个模型学习，既可以把多个来源的行为转入承接模型，也可以约束更新后系统不要丢失旧行为。\term{经验回放}（experience replay）在后续训练中重新使用保存的历史经验，使旧要求继续出现在学习依据中。\term{参数约束}（parameter constraint）限制更新偏离参考参数的范围，表达对某些既有计算的保留要求。\term{模块路由}（module routing）根据任务或输入选择参与计算的模块，使不同请求可以采用不同的共享范围。操作名称没有规定最终研究目标，也不自动保证保留或迁移成功。

例如，“撤回一张历史训练照片的影响”主要调用更新与维护以及可信性部分的删除标准，而不是只删除缓存文件；“把两个地区模型放进同一部署预算”主要调用整合，并回到来源选择与成本评价；“每月加入两种新花，同时维持已有类别”则同时调用目标适应与长期维护。若请求涉及动作、奖励和环境交互，探索及策略的专门推导由第~\ref{part:reinforcement-learning}部分承担；存储、通信、分布式训练与上线实现由第~\ref{part:systems}部分承接。本部分关心的是在这些条件下，已有知识究竟怎样发挥作用。

\section{知识的演进与迁移的研究历史}
\label{sec:ket-overview-history}

知识复用并不是在大规模预训练出现后才产生的问题。共享表示、学习偏置、顺序干扰和多模型组合有不同的研究入口，后来又在同一系统中交汇。下面沿四条线索观察代表工作分别扩大了什么问题范围；同一线索内部按时间展开，工作之间不构成相互替代的阶段，也不据选取的论文认定某个思想的唯一来源。

\subsection{共享规律与可迁移表示的研究}
\label{sec:ket-overview-history-representation}

单个任务的样本不足时，其他任务的训练要求能够提供额外限制。Caruana在1997年的多任务学习工作中，系统讨论通过共同表示让相关任务的训练信号影响彼此的学习。放在花卉例子中，物种和病害标签虽然回答不同问题，却可能共同要求网络保留叶片或花瓣的某些细节；共享层使这些要求作用于同一个表示。这里的收益来自任务间可以利用的联系，不是“输出头越多越好”，也不是把独立训练的模型简单放在一起\scite{caruana1997}

Baxter在2000年的归纳偏置学习研究中，把问题进一步写成从相关任务环境中选择适合后续学习的假设空间。\term{归纳偏置}（inductive bias）是学习者在多个可能解释之间作出选择的倾向，例如偏好某种特征表示或某类函数。若能用多个任务判断哪些表示共同有用，新任务就不必从所有可能表示中重新搜索。但这一论证依赖任务来源的关系，以及候选假设空间族的复杂度限制；它不是仅由训练任务数量增加就得到的普遍结论。其重要区分是：在当前几个任务上表现良好，与为同一环境中的未见任务准备合适的学习空间，并非同一个评价目标\scite{ket-baxter2000-bias}

当来源和目标不仅任务不同，数据分布也不同，来源成绩便不足以代表目标收益。2010年的领域自适应不可学习性研究明确分析了输入分布关系、共同低误差预测器和条件标签假设各自能提供什么。即使无标签输入看起来接近，算法仍可能无法辨认目标标签规律；反之，输入变化也不必然使一切来源信息失效。这条研究线索要求把表示的通用性与适用条件一起讨论，而不是只寻找一个不受任何环境影响的特征空间\scite{bendavid2010impossibility}

Yosinski等人在2014年的实验研究把可迁移性落实到深度网络的具体层及其使用方式。在其自然图像卷积网络设置中，较早层与较后层呈现不同程度的任务特异性；迁移受损既可能来自特征偏向原任务，也可能来自切开原网络中协同形成的计算关系。因而，“取出中间层”不是一个没有代价的操作，同一层冻结使用和继续微调也回答不同问题。该研究支持按层、按目标和按使用方式检查表示，而不能据此规定所有架构中某个固定层数都通用\scite{yosinski2014}

这些工作把“共享什么”从笼统的任务相关性逐步落实为表示、假设空间和参数的可用范围。遇到分布差异时，还需要准备什么对应依据、怎样对齐或重加权，交由第~\ref{chap:ket-update-maintenance}章展开；本条线索保留的判断是，通用性始终相对于目标要求和可用信息而言。

\subsection{快速适应与学习过程复用的研究}
\label{sec:ket-overview-history-adaptation}

一个表示在大量目标标签下能学好，并不意味着几张新物种照片就足以确定分类边界。快速适应研究由此把关注点从已有预测器转向准备阶段：应当预先学好比较样本的方式，找到适合少量更新的参数起点，还是直接学习怎样更新？三者分别改变目标学习中最需要补充的内容，不能仅凭最终都用了少量样本就视为同一机制。

2016年Andrychowicz等人的学习优化器研究把更新规则本身作为可学习对象：让一个参数化更新器从训练问题中学习如何利用梯度等信息，而不只依赖手工给定的优化规则。这里保留下来的是产生参数变化的规则，其泛化需要在新的优化问题上检验。一个更新器在训练所覆盖的问题结构上有效，不意味着它对任意网络规模、损失或训练时长都适用\scite{ket-andrychowicz2016-learning}

2017年的MAML，即\term{模型无关元学习}（model-agnostic meta-learning），选择另一种准备内容：训练一个经过少量目标梯度更新后能够有效预测的初始化。它不是只让初始模型在混合任务样本上误差小，而是让适应后的表现反过来决定起点。同年的\term{原型网络}（prototypical networks）则学习适于比较的表示，再用少量支持样本形成各类别的代表位置。前者把少量参数更新纳入训练目标，后者让表示中的距离支持少样本分类；两者说明准备内容可以是适应起点，也可以是便于目标端估计的度量空间\scite{finn2017,ket-snell2017-prototypical}

目标模型规模增大后，快速适应还遇到每个任务要保存和更新多少参数的问题。Houlsby等人在2019年的工作中使用\term{适配器}（adapter），即加入原网络、由目标任务训练的小型模块，在保持原网络参数固定的条件下提供任务特有变化。这类\term{参数高效适应}（parameter-efficient adaptation）关心用较少可训练参数承载目标调整；它不自动意味着标签更少、训练时间更短或推理成本更低，这些量仍需分别核算\scite{ket-houlsby2019-adapters}

2020年的GPT-3研究展示了另一种目标端使用方式：在所报告的少样本设置中，任务说明和示例通过文本进入输入，预测时不进行梯度更新。这把“给模型少量经验”与“用这些经验修改权重”分开了，也使上下文组织成为适应的一部分\scite{brown2020}

在广泛数据上预训练、能够为多种下游任务提供共同起点的\term{基础模型}（foundation model），进一步扩大了这些准备方式的使用范围，但没有消除它们的区别。度量表示、初始化与更新规则是准备侧的不同对象；附加参数训练与上下文条件化是目标侧的不同操作。少样本表现良好，既可能主要得益于来源表示已经适合目标，也可能来自有效使用支持样本，不能单凭一个结果认定模型学会了新的学习算法。第~\ref{chap:ket-transfer-adaptation}章据此分开解释准备、目标适应及其对照。

\subsection{顺序学习与长期知识维护的研究}
\label{sec:ket-overview-history-maintenance}

当同一组参数需要先后承担不同要求，新更新可能破坏旧计算。\term{灾难性遗忘}（catastrophic forgetting）指后续学习引起既有能力的大幅丧失，而不只是记忆随时间缓慢减弱。French在1999年的综述中已经系统讨论这一现象及表示重叠、学习干扰和解决途径。顺序学习中的保留需求因而早于后来深度表示迁移的普及；它与共享规律研究长期并行，而不是后者完成以后才出现的新阶段\scite{ket-french1999-forgetting}

2017年发表的\term{弹性权重巩固}（elastic weight consolidation, EWC）把保留要求落实为对参数变化的差别约束：对过去任务较重要的参数，更新时施加较强限制，对其他参数留下更多调整空间。它使用历史参数及其重要性估计表达旧要求，不等于保存所有旧样本，也不等于确认某个参数只对应某条事实。原论文的预印本发表于2016年，本文沿用正式发表年份2017年；其实验说明了相应设置中的保留效果，不构成任意长任务序列上不遗忘的保证\scite{kirkpatrick2017ewc}

参数约束之外，经验回放使旧样本或其替代信息重新进入训练，结构隔离则通过任务分支、专属参数等方式减少直接覆盖。上述参数约束依赖重要性估计，回放依赖可保存的历史经验，隔离依赖可分配的容量。约束过强可能妨碍新学习，回放受历史覆盖和预算限制，隔离则增加保存与选择模块的负担。长期维护还需要检查可塑性：系统当前仍答得对，不代表它以后还能以相同效率学会新的要求。保留成绩、吸收当前信息和维持后续学习能力因而不能合成一个“训练更久”的指标。

显式知识系统提供了不同于参数保留的视角。Mitchell等人在2018年的《Never-Ending Learning》中介绍\term{永不停止的语言学习系统}（Never-Ending Language Learner, NELL）：它从网络文本中积累类别与关系信念，记录置信度及出处，并利用已学关系支持后续抽取和推理。论文报告该系统自2010年开始长期运行，2018年是这里所引论文的发表时间，不是系统开始学习的时间。该案例把事实记录的增长、错误修正和学习机制改进放在同一个长期过程中考察，而不只观察一次训练后的分类准确率\scite{ket-mitchell2018-nell}

在持续预训练与连续使用中，保留要求还会随事实有效性而变化。\term{知识编辑}（knowledge editing）针对指定事实或行为提出定向修改，例如示例~\ref{ex:ket-overview-dated-fact}中的负责人交接；它应保留历史时间条件下仍成立的回答，而不是把甲在所有问题中的出现都抹掉。\term{机器遗忘}（machine unlearning）则要求按规定标准消除指定训练记录的影响，不等同于把一条事实改成另一个答案。由此，“保持旧知识”需要更精确地表述为“保持当前仍然有效且应当保留的知识”，同时允许纠错、撤回与必要的关联变化。

2024年的持续学习研究把注意力从“旧任务是否还答得对”推进到\term{可塑性损失}：模型即使保持已有成绩，也可能越来越难从后续数据学习。Dohare等在长序列设置中系统展示了这种现象及持续反向传播等干预，提醒维护评价必须把保留与后续学习速度并列；其结论仍受所测网络、任务流和优化设置限制\scite{ket-dohare2024-plasticity}

这条线索连接参数、行为和外部记录，却不要求用一种技术解决全部维护问题。知识图谱存储和本体版本管理负责显式记录的专门结构，本部分只保留它们与学习、检索和版本验证的接口。第~\ref{chap:ket-update-maintenance}章具体讨论更新范围、关联一致性与多次修改后的维护，第~\ref{sec:privacy-unlearning-validation}节给出训练影响消除不能省略的验证边界。

\subsection{多来源能力整合的研究}
\label{sec:ket-overview-history-integration}

多个分别训练的模型已经存在时，可以在预测阶段汇集它们的判断。\term{预测集成}（prediction ensemble）保留这些来源的计算，再根据明确规则组合输出。它利用的是来源误差或覆盖的互补性，但来源数量增加并不自动创造互补信息：对同一背景产生相同错误的两个地区模型，未必能靠平均解决问题。保留全部来源还意味着推理时需要承担相应执行与存储成本。

如果希望部署时只保留一个模型，就需要把来源行为转入另一载体。Hinton、Vinyals与Dean在2015年的蒸馏工作中进一步发展了先前模型压缩的思路，用教师输出中的相对概率指导学生学习。与只给出最大概率类别相比，软化的输出保留了教师对其他类别的相对倾向。学生通过这些监督信号学习，并不需要与教师具有相同的参数坐标；但查询输入覆盖不足或学生容量受限时，不能保证保留所有来源能力。此处的2015年采用所引arXiv版本的公开年份\scite{hinton2015distilling}

整合也可以发生在中间计算层面。Alet等人在2018年的模块化元学习研究中学习可供重组的神经模块，再通过新的组合适应新的任务。这连接了两条问题线索：准备模块库时需要利用任务经验，面对目标时则需要在允许的结构中选择与连接现成模块。它不是把任意两个训练模型接起来就能工作；接口、组合语法以及目标证据决定可搜索的范围，论文中的组合收益也有具体任务条件\scite{ket-alet2018-modular}

例如，花卉系统已有识别颜色、叶片形状与位置的部件，新的请求是“找出图中具有锯齿叶缘的白色花朵”。即使各属性在单独测试中识别良好，仍要检查系统能否把属性绑定到同一对象，而不是把一个对象的颜色与另一个对象的叶缘混在一起。未见属性组合、指令组合与技能重组的收益，需要针对组合本身设立证据；它们不能由来源任务准确率相加推得。

2022年的Model Soups研究进一步考察直接组合参数：在同一预训练起点的多个微调结果之间平均权重，从而形成单个模型。其报告的收益建立在特定模型、任务及微调解的条件上，不能推广为任意独立模型的参数平均都有效。参数平均后的单模型预测通常也不同于多个模型预测的平均，前者在网络计算前组合参数，后者在各来源完成计算后组合输出。共同起点有助于形成可比较的参数坐标，但不是来源能力全部保留的充分条件\scite{ket-wortsman2022-soups}

基础模型及其适配器库扩大了可获得的来源集合，使中间模块、完整参数及相对底座的参数增量都成为整合对象。保留原来源共同运行，需要持续访问各来源；形成承接模型，则需要在准备时获得足够的来源行为、参数或表示，并检查转入后的损失。无论采用哪条路径，来源能力保持、任务冲突、未见组合表现和全过程成本都要分别评价。第~\ref{chap:ket-integration}章集中解释这些机制；来源规模增长只是改变可用条件，不能直接写成知识整合已经解决。

到2025年，低秩适配器的整合进一步暴露出因子坐标不一致问题。KnOTS先在完整增量层面建立共同子空间，再在该坐标中应用融合规则，说明“参数形状相同”仍不足以保证低秩知识能够直接相加\scite{ket-stoica2025knots}这类工作扩展了整合对象，却没有消除独立来源基准、异质模型和组合能力评价的需要。

图~\ref{fig:ket-overview-history}把上述四条线索置于共同时间轴。点的位置对应已核验的代表年份，带内没有表示方法替代的箭头；跨带虚线只说明共享的问题，不表示历史上的直接继承关系。例如，学到可复用的表示与准备快速适应的起点，都依赖任务之间保留什么规律；参数保留与多来源整合则都要检查既有能力是否在改变之后仍然可用。

\begin{figure*}[htbp]
  \centering
  % Equal year scale: x = 36 + 4 * (year - 1997), in mm.
% Event years are publication years except the explicitly dated arXiv distillation.
% Bands show concurrent research questions, not their starting/ending dates.
% Dashed arrows link questions, not claims of historical influence.
\begin{tikzpicture}[
  x=1mm,y=1mm,
  font=\figurefont\footnotesize,
  every node/.style={text=BookInk},
  band title/.style={
    align=center,text width=27mm,minimum height=10mm,
    inner sep=0.5mm,font=\figurefont\small
  },
  event/.style={align=center,inner sep=0.7mm},
  leader/.style={draw=BookMuted,line width=0.6pt},
  relation/.style={
    draw=BookTeal,line width=0.8pt,dashed,
    -{Latex[length=1.8mm,width=1.3mm]}
  }
]
  \foreach \level in {0,-32,-64,-96} {
    \fill[BookPaper] (33,\level-12) rectangle (162,\level+12);
    \draw[BookRule,line width=0.6pt] (35,\level) -- (158,\level);
  }
  \draw[BookMuted,line width=0.7pt] (36,19) -- (150,19);
  \foreach \year in {1997,2000,2005,2010,2015,2020,2025} {
    \draw[BookMuted,line width=0.6pt]
      ({36+4*(\year-1997)},18) -- ({36+4*(\year-1997)},20);
    \node[anchor=south] at ({36+4*(\year-1997)},20) {\year};
  }
  \node[anchor=east,text=BookMuted] at (29,21) {年份};

  \node[band title] (representation) at (14,0) {共享规律\\与可迁移表示};
  \node[band title] (adaptation) at (14,-32) {快速适应\\与学习过程复用};
  \node[band title] (maintenance) at (14,-64) {顺序学习\\与长期知识维护};
  \node[band title] (integration) at (14,-96) {多来源\\能力整合};

  \draw[relation] (representation.south) -- (adaptation.north);
  \node[anchor=west,text=BookTeal] at (15,-16) {任务规律};
  \draw[relation] (maintenance.south) -- (integration.north);
  \node[anchor=west,text=BookTeal] at (15,-80) {能力保持};

  \foreach \name/\year/\level/\ink in {
    caruana/1997/0/BookBlue,
    baxter/2000/0/BookBlue,
    yosinski/2014/0/BookBlue,
    maml/2017/-32/BookTeal,
    adapters/2019/-32/BookTeal,
    gpt/2020/-32/BookTeal,
    french/1999/-64/BookAmber,
    ewc/2017/-64/BookAmber,
    nell/2018/-64/BookAmber,
    plasticity/2024/-64/BookAmber,
    distillation/2015/-96/BookBlue,
    modular/2018/-96/BookBlue,
    soups/2022/-96/BookBlue,
    knots/2025/-96/BookBlue%
  } {
    \coordinate (\name) at ({36+4*(\year-1997)},\level);
    \fill[\ink] (\name) circle[radius=1mm];
  }

  \node[event,anchor=south] (caruana label) at (47,4) {多任务学习\\Caruana（1997）};
  \draw[leader] (caruana) -- (36,2) -| (caruana label.south);
  \node[event,anchor=north] (baxter label) at (67,-4) {偏置学习\\Baxter（2000）};
  \draw[leader] (baxter) -- (49.2,-2) -| (baxter label.north);
  \node[event,anchor=south] (yosinski label) at (118,4) {特征可迁移性\\Yosinski等（2014）};
  \draw[leader] (yosinski) -- (110.8,2) -| (yosinski label.south);

  \node[event,anchor=south] (maml label) at (108,-28) {MAML\\2017};
  \draw[leader] (maml) -- (124,-30) -| (maml label.south);
  \node[event,anchor=north] (adapters label) at (114,-36) {适配器\\2019};
  \draw[leader] (adapters) -- (132.8,-34) -| (adapters label.north);
  \node[event,anchor=south] (gpt label) at (149,-28) {GPT-3上下文\\2020};
  \draw[leader] (gpt) -- (137.2,-30) -| (gpt label.south);

  \node[event,anchor=north] (french label) at (64,-68) {遗忘研究回顾\\French（1999）};
  \draw[leader] (french) -- (44.8,-66) -| (french label.north);
  \node[event,anchor=south] (ewc label) at (109,-60) {EWC\\2017};
  \draw[leader] (ewc) -- (124,-62) -| (ewc label.south);
  \node[event,anchor=north] (nell label) at (144,-68) {NELL案例论文\\2018};
  \draw[leader] (nell) -- (120,-66) -| (nell label.north);
  \node[event,anchor=south] (plasticity label) at (145,-60) {可塑性损失\\2024};
  \draw[leader] (plasticity) -- (144,-62) -| (plasticity label.south);

  \node[event,anchor=south] (distillation label) at (100,-92) {知识蒸馏\\2015};
  \draw[leader] (distillation) -- (108,-94) -| (distillation label.south);
  \node[event,anchor=north] (modular label) at (118,-100) {模块化元学习\\2018};
  \draw[leader] (modular) -- (120,-98) -| (modular label.north);
  \node[event,anchor=south] (soups label) at (137,-92) {Model Soups\\2022};
  \draw[leader] (soups) -- (136,-94) -| (soups label.south);
  \node[event,anchor=north] (knots label) at (149,-100) {KnOTS\\2025};
  \draw[leader] (knots) -- (148,-98) -| (knots label.north);
\end{tikzpicture}

  \caption{知识演进与迁移的四条研究线索。横轴使用等距年份，圆点为正文讨论的代表工作，标签用引线避让；横带表示问题长期并行，不表示研究的起止年份。跨带虚线的文字说明共同问题，而非断言论文的直接影响关系。EWC采用2017年正式发表年份，NELL采用2018年案例论文年份，蒸馏采用2015年arXiv版本年份；图不声明任何工作的“首次”地位。}
  \label{fig:ket-overview-history}
\end{figure*}

\section{本章小结}
\label{sec:ket-overview-summary}

花卉系统进入新地区时，需要检查原预测规律能否继续使用；增加病害或定位任务时，需要重新明确输出语义；连续加入新物种时，还要在形成新能力的同时保持累积能力。负责人交接则要求按有效日期修正事实，不必由数据分布变化触发。它们共同说明，决定研究问题的是当前要实现和保留什么，而不是是否又运行了一轮训练。

知识载体、任务能力与语义事实不能混同。样本、表示、参数、模块、教师行为和学习规则可以提供不同形式的经验，但其价值须相对于目标和使用方式判断。来源与目标说明经验从哪里来、要服务什么要求，信息条件说明算法实际能看到和改变什么。临时使用上下文、更新外部记忆与修改参数可以共同改变系统行为，却需要不同的持久性和有效范围证据。

后续章节在这些共同设定下判断和选择复用内容，研究更新、整合、迁移与共享，并统一评价目标收益、保留要求及全过程代价。同一种算法可以服务多个研究目标，同一个系统也可以反复经历这些问题。只有把目标、来源关系和信息权限写清楚，才能确定合理对照，并判断已有经验究竟帮助了什么、又在哪些条件下失效。
